\section{Обсуждение результатов}
\label{sec:discussion}

\subsection{Роль квазиньютоновских методов}

BFGS и SciPy BFGS работают с полным градиентом и используют приближение обратного гессиана, то есть частично учитывают локальную кривизну функции потерь. Для обучения потенциалов это особенно важно: функционал одновременно содержит энергии и силы, а поверхность loss может быть плохо обусловленной. На MTP это даёт наиболее заметное преимущество BFGS-группы над mini-batch методами. В рассматриваемой конфигурации на AgPd минимальную validation EPA RMSE получил native BFGS, поэтому для MTP основной вывод достаточно устойчив: full-batch квазиньютоновские методы остаются сильнее first-order mini-batch оптимизаторов.

Для ETN native BFGS также сохраняет роль сильного baseline. При этом результат SciPy BFGS уже не повторяет поведение MTP: внешний BFGS-интерфейс не всегда улучшает качество и может давать большее время обучения при близкой или худшей validation error. Поэтому выводы о BFGS-вариантах нельзя переносить между MTP и ETN без отдельной проверки.

\subsection{Поведение SGD, Momentum и Adam}

SGD и Momentum в проведённых экспериментах сильно зависят от learning rate, batch size и clipping. Шум mini-batch градиента затрудняет движение по поверхности loss, а слишком крупный шаг может приводить к плато или нестабильности. Momentum сглаживает направление, но не отменяет необходимости аккуратно подбирать масштаб шага.

Adam в большинстве mini-batch сравнений оказывается сильнее SGD и Momentum. Это согласуется с его устройством: метод использует сглаженный первый момент и адаптивную нормировку по второму моменту. Для MTP на AgPd Adam даёт минимальную validation EPA RMSE среди mini-batch методов. Для ETN на AgPd он также лидирует по validation EPA RMSE, но по validation forces RMSE меньшую ошибку дают варианты Muon.

\subsection{Различие MTP и ETN}

MTP опирается на аналитически заданный набор инвариантных базисных функций. Такая структура хорошо сочетается с full-batch квазиньютоновскими методами и делает масштаб инициализации особенно заметным для SciPy BFGS. Для MTP native BFGS и SciPy BFGS близки по energy error, но native BFGS лучше по validation forces RMSE. ETN устроен иначе: он использует тензорно-сетевую параметризацию эквивариантных признаков. Это повышает гибкость модели, но меняет ландшафт оптимизации и относительное поведение методов.

Практически это означает, что оптимизатор нельзя выбирать отдельно от архитектуры потенциала. Это видно не только на SciPy BFGS, который хорошо работает для MTP при подходящей инициализации, но не становится универсально лучшим методом для ETN. Ещё более показательно ведёт себя Muon. Для MTP на AgPd варианты Muon уступают Adam и BFGS-вариантам по validation EPA RMSE. Для ETN картина меняется: на AgPd Muon почти сравнивается с Adam по energy error и превосходит его по force error, а на MoNbTaVW по validation EPA RMSE превосходит Adam и BFGS-варианты.

Такое различие хорошо согласуется со структурой моделей. Muon нормализует матричные направления обновления, поэтому для ETN с тензорно-сетевой параметризацией и большим числом матричных блоков такой тип шага выглядит более естественным, чем для MTP, где параметры связаны с явно заданными базисными инвариантами. Поэтому вывод о лучшем mini-batch методе нельзя сводить только к Adam. Adam остаётся сильным baseline, но для ETN варианты Muon становятся конкурентными, а на MoNbTaVW дают меньший средний validation EPA RMSE, чем все остальные рассмотренные методы. Даже с учётом разброса ансамбль Muon на MoNbTaVW расположен в области меньших energy errors, чем ансамбль Adam.

\subsection{Перенос выводов на MoNbTaVW}

MoNbTaVW нужен как проверка выводов на другом, более многокомпонентном датасете. В ETN-экспериментах native BFGS и SciPy BFGS остаются сильными full-batch baseline, особенно по validation forces RMSE и времени обучения. Однако после \(10000\) mini-batch шагов разница между full-batch и mini-batch подходами по energy error не только уменьшается: варианты Muon дают меньшую validation EPA RMSE, чем native BFGS и SciPy BFGS, хотя всё ещё уступают им по force error и времени обучения.

Итог по mini-batch методам становится датасет-зависимым и зависит от целевой метрики. На AgPd Adam лидирует по validation EPA RMSE, а Muon даёт меньшую force error для ETN. На MoNbTaVW варианты Muon pad sqrt и Muon factorization существенно лучше Adam, SGD и Momentum по validation EPA RMSE и становятся лучшими по energy error среди всех рассмотренных методов. Поэтому Muon нельзя рассматривать как второстепенный метод: для ETN он показывает выраженный потенциал, особенно на MoNbTaVW. При этом по force error и времени обучения BFGS-варианты пока остаются сильнее.
